\subsection{直态射}

定义 1. --- 设 E 是拓扑向量空间。E 的向量子空间称为直的，如果它有一个拓扑补。

为使 E 的向量子空间 F 是直的，必须且只需在 \(\mathcal{L}(\mathrm{E})\) 中存在一个以 F 为像的投影。如果 \(p\) 是这样一个投影，E 是 F 与 \(p\) 的核的拓扑直和。\(1_{\mathrm{E}} - p\) 的核是 F；因此如果 E 是分离的，那么 F 是闭的。

设 F 是 E 的直向量子空间，\(E_{1}\) 是 E 的包含 F 的向量子空间；那么 F 是 \(E_{1}\) 的直向量子空间。

E 中每个有有限余维的闭向量子空间都是直和项（EVT, I, p. 15, prop. 3），而 E 的每个包含在 \{0\} 的闭包中的向量子空间 F 都是直的（事实上，存在一个以 F 为像的投影，且这样一个投影必然是连续的）。

命题 1. --- 设 E 是局部凸空间。E 的每个有限维向量子空间都有补。

设 F 是 E 的有限维向量子空间。用 S 表示 \(\{0\}\) 在 F 中的闭包的一个拓扑补（该闭包是 E 的直子空间）。空间 \(F_{2} = F \cap S\) 是分离的且有限维；因此存在整数 \(n \geqslant 0\) 使得 \(F_{2}\) 同构于 \(K^{n}\)（EVT, I, p. 14, th. 2）。于是存在一个连续线性映射 \(p: S \to F_{2}\)，它延拓 \(F_{2}\) 的恒等映射。p 的核是 \(F_{2}\) 在 S 中的拓扑补，也是 F 在 E 中的拓扑补。

定义 2. --- 设 E 与 F 是局部凸空间，u 是从 E 到 F 中的连续线性映射。称 u 是直态射，如果 u 是严格态射，其核是 E 的直向量子空间且其像是 F 的直向量子空间。

设 \(u \in \mathcal{L}(\mathrm{E};\mathrm{F})\)。为使 \(u\) 的核与像分别是 E 与 F 的直向量子空间，必须且只需存在分解为拓扑直和的 \(\mathrm{E} = \mathrm{E}_1 \oplus \mathrm{E}_2\) 与 \(\mathrm{F} = \mathrm{F}_1 \oplus \mathrm{F}_2\)，使得 \(u\) 由矩阵 \(\begin{pmatrix} u_1 & 0 \\ 0 & 0 \end{pmatrix}\) 表示，其中 \(u_1 \in \mathcal{L}(\mathrm{E}_1; \mathrm{F}_1)\) 是双射。由于这时 \(u\) 的核是 \(\mathrm{E}_2\)，从 \(\mathrm{E}_1\) 到 \(\mathrm{E}/\mathrm{E}_2\) 上的典范线性映射是同构，且 \(u\) 的像是 \(\mathrm{F}_1\)，可以看出 \(u\) 是严格态射当且仅当 \(u_1\) 是从 \(\mathrm{E}_1\) 到 \(\mathrm{F}_1\) 上的同构。用 \(v\) 表示 \(\mathcal{L}(\mathrm{F};\mathrm{E})\) 中由矩阵 \(\begin{pmatrix} u_1^{-1} & 0 \\ 0 & 0 \end{pmatrix}\) 表示的元素。那么线性映射 \(u \circ v\) 是 F 中以 \(\mathrm{F}_2\) 为核、以 \(\mathrm{F}_1\) 为像的投影，线性映射 \(v \circ u\) 是 E 中以 \(\mathrm{E}_2\) 为核、以 \(\mathrm{E}_1\) 为像的投影，且有 \(u \circ v \circ u = u\)。

反之，我们有下面的结果：

命题 2. --- 设 E 与 F 是局部凸空间，设 \(u \in \mathcal{L}(\mathrm{E};\mathrm{F})\) 与 \(v \in \mathcal{L}(\mathrm{F};\mathrm{E})\)。假设 \(u = u \circ v \circ u\)。那么 \(u\) 是直态射。此外，\(v \circ u\) 是 E 中以 \(\operatorname{Ker}(u)\) 为核的连续投影，而 \(u \circ v\) 是 F 中以 \(\operatorname{Im}(u)\) 为像的连续投影。

置 \(p = v \circ u\) 与 \(q = u \circ v\)。我们有

\[
p ^ {2} = v \circ (u \circ v \circ u) = v \circ u = p, q ^ {2} = (u \circ v \circ u) \circ v = u \circ v = q,
\]

因而 p 与 q 是投影。它们是连续的。设 \(E_{1}\) 与 \(E_{2}\) 是 p 的像与核，设 \(F_{1}\) 与 \(F_{2}\) 是 q 的像与核。由于 \(u \circ v \circ u = u\)，我们有

\[
\operatorname{Ker} (u) \subset \operatorname{Ker} (v \circ u) \subset \operatorname{Ker} (u \circ v \circ u) = \operatorname{Ker} (u),
\]

因而 \(\operatorname{Ker}(u) = \operatorname{Ker}(v \circ u) = \operatorname{Ker}(p) = \operatorname{E}_{2}\)。类似地，由包含关系

\[
\operatorname{Im} (u) \supset \operatorname{Im} (u \circ v) \supset \operatorname{Im} (u \circ v \circ u) = \operatorname{Im} (u)
\]

得出 \(u\) 的像是 \(\mathrm{F}_{1}\)。空间 \(\mathrm{E}_{2}\) 与 \(\mathrm{F}_{1}\) 分别是 E 与 F 的直向量子空间。

向量空间 E 是 \(E_{1}\) 与 \(E_{2}\) 的拓扑直和，且 u 通过限制定义了一个双射线性映射 \(u_{1}\)（它从 \(E_{1}\) 到 \(F_{1}\) 上）。对每个 \(x \in E_{1}\)，我们有 \(v(u(x)) = p(x) = x\)，且映射 \(u_{1}^{-1}\)（它从 \(F_{1}\) 到 \(E_{1}\) 上）在 \(F_{1}\) 上与 v 一致。因此，\(u_{1}\) 是从 \(E_{1}\) 到 \(F_{1}\) 上的同构，而 u 是从 E 到 F 中的严格态射。

设 E 与 F 是局部凸空间，且 \(u \in \mathcal{L}(\mathrm{E};\mathrm{F})\)。为使 u 是单射（相应地，满射）的直态射，必须且只需存在一个 \(v\)（它在 \(\mathcal{L}(\mathrm{F};\mathrm{E})\) 中），使得 \(v \circ u = 1_{E}\)（相应地，\(u \circ v = 1_{F}\)）（参见 TG, III, p. 47 及 48）。

命题 3. --- 设 E 与 F 是巴拿赫空间。下列集合都是巴拿赫空间 \(\mathcal{L}(\mathrm{E};\mathrm{F})\) 的开子集：

\begin{enumerate}
\def\labelenumi{\alph{enumi})}
\item
  从 E 到 F 上的同构组成的集 \(\mathcal{I}(\mathrm{E};\mathrm{F})\)；
\item
   从 E 到 F 中的单射直态射组成的集 \(\mathcal{M}\mathcal{D}(\mathrm{E};\mathrm{F})\)，且更精确地说，对 F 的每个闭向量子空间 \(F_{1}\)，集 \(\mathcal{M}_{\mathrm{F}_{1}}(\mathrm{E};\mathrm{F})\)（它由 \(\mathcal{M}\mathcal{D}(\mathrm{E};\mathrm{F})\) 中像为 \(F_{1}\) 的拓扑补的元素组成）；
\item
   从 E 到 F 中的满射直态射组成的集 \(\mathcal{E}\mathcal{D}(\mathrm{E};\mathrm{F})\)，且更精确地说，对 E 的每个闭向量子空间 \(\mathrm{E}_1\)，集 \(\mathcal{E}_{\mathrm{E}_1}(\mathrm{E};\mathrm{F})\)（它由 \(\mathcal{E}\mathcal{D}(\mathrm{E};\mathrm{F})\) 中核为 \(\mathrm{E}_1\) 的拓扑补的元素组成）。
\end{enumerate}

此外，映射 \(u \mapsto u^{-1}\)（它从 \(\mathcal{I}(\mathrm{E};\mathrm{F})\) 到 \(\mathcal{I}(\mathrm{F};\mathrm{E})\) 上）是解析的。

依定义，\(\mathcal{I}(\mathrm{E};\mathrm{E})\) 是完备赋范代数 \(\mathcal{L}(\mathrm{E})\) 的可逆元素组成的集。由 TG, IX, p. 40, prop. 14，它是 \(\mathcal{L}(\mathrm{E})\) 的开子集。映射 \(u\mapsto u^{-1}\)（它从 \(\mathcal{I}(\mathrm{E};\mathrm{E})\) 到 \(\mathcal{I}(\mathrm{E};\mathrm{E})\) 中）是解析的（LIE, III, § 1, 第 1 目, 例子 2）。

如果集 \(\mathcal{I}(\mathrm{E};\mathrm{F})\) 是空的，那么它是开的。否则，设 \(u_0\) 是从 E 到 F 上的一个同构。于是映射 \(v\mapsto u_0\circ v\) 是从 \(\mathcal{L}(\mathrm{E})\) 到 \(\mathcal{L}(\mathrm{E};\mathrm{F})\) 上的同构，它把 \(\mathcal{I}(\mathrm{E};\mathrm{E})\) 变换为 \(\mathcal{I}(\mathrm{E};\mathrm{F})\)。因此集 \(\mathcal{I}(\mathrm{E};\mathrm{F})\) 在 \(\mathcal{L}(\mathrm{E};\mathrm{F})\) 中是开的。此外，如果 \(u = u_0\circ v\) 是 \(\mathcal{I}(\mathrm{E};\mathrm{F})\) 的一个元素，我们有 \(u^{-1} = v^{-1}\circ u_0^{-1}\)，因而映射 \(u\mapsto u^{-1}\)（它从 \(\mathcal{I}(\mathrm{E};\mathrm{F})\) 到 \(\mathcal{I}(\mathrm{F};\mathrm{E})\) 上）是解析的。

设 \(F_{1}\) 是 F 的闭向量子空间。对每个 \(u \in \mathcal{L}(E;F)\)，用 \(\overline{u}\) 表示 \(\mathcal{L}(E \times F_{1};F)\) 中由 \(\overline{u}(x,y) = u(x) + y\) 定义的元素。映射 \(u \mapsto \overline{u}\)（它从 \(\mathcal{L}(E;F)\) 到 \(\mathcal{L}(E \times F_{1};F)\) 中）是连续的，而 \(\mathcal{M}_{\mathrm{F}_{1}}(E;F)\) 是由 \(\mathcal{L}(E;F)\) 中使 \(\overline{u}\) 属于 \(\mathcal{I}(E \times F_{1};F)\) 的元素 u 组成的集。由于 \(\mathcal{I}(E \times F_{1};F)\) 在 \(\mathcal{L}(E \times F_{1};F)\) 中是开的

由前所述，集 \(\mathcal{M}_{\mathrm{F}_1}(\mathrm{E};\mathrm{F})\) 在 \(\mathcal{L}(\mathrm{E};\mathrm{F})\) 中是开的。对 \(\mathcal{MD}(\mathrm{E};\mathrm{F})\) 同样如此，它是诸 \(\mathcal{M}_{\mathrm{F}_1}(\mathrm{E};\mathrm{F})\)（当 \(\mathrm{F}_1\) 取遍 F 的闭向量子空间组成的集时）的并。

设 \(E_{1}\) 是 E 的闭向量子空间，p 是从 E 到商巴拿赫空间 \(E/E_{1}\) 上的典范映射。为使一个元素 \(u\)（它在 \(\mathcal{L}(E;F)\) 中）属于 \(\mathcal{E}_{\mathrm{E}_{1}}(\mathrm{E};\mathrm{F})\)，必须且只需映射 \((u,p)\)（它从 E 到 \(F\times E/E_{1}\) 中）属于 \(\mathcal{I}(E;F\times E/E_{1})\)。与前面一样，我们推出 \(\mathcal{E}_{\mathrm{E}_{1}}(\mathrm{E};\mathrm{F})\) 在 \(\mathcal{L}(\mathrm{E};\mathrm{F})\) 中是开的；对 \(\mathcal{E}\mathcal{D}(\mathrm{E};\mathrm{F})\) 同样如此，它是诸 \(\mathcal{E}_{\mathrm{E}_{1}}(\mathrm{E};\mathrm{F})\) 的并。

